\documentclass[12pt]{article}
\usepackage[utf8]{inputenc}
\usepackage[T1]{fontenc}
\usepackage[french]{babel}
\usepackage{amsmath}
\usepackage[version=4]{mhchem}
\usepackage{graphicx}
\usepackage{float}
\usepackage{hyperref}

\title{Résumé : Réactions Acido-Basiques en Solution Aqueuse}
\author{}
\date{}

\begin{document}

\maketitle

\section{Introduction}
Les réactions acido-basiques sont des réactions chimiques qui impliquent un transfert de protons (\ce{H+}).

\section{Acides et Bases}
\subsection{Définitions}
\begin{itemize}
    \item Un \textbf{acide} est une espèce chimique capable de libérer un ion \ce{H+}.
    \item Une \textbf{base} est une espèce chimique capable de capter un ion \ce{H+}.
\end{itemize}

\subsection{Demi-équations}
Pour les réactions acido-basiques, les demi-équations sont souvent utilisées pour décrire les transformations :
\begin{itemize}
    \item Pour un acide : \ce{AH <=> A- + H+}
    \item Pour une base : \ce{B + H+ <=> BH+}
\end{itemize}

\subsection{Exemples}
\begin{itemize}
    \item Acide acétique : \ce{CH3COOH <=> CH3COO- + H+}
    \item Ammoniac : \ce{NH3 + H+ <=> NH4+}
\end{itemize}

\subsection{Couples Acide/Base}
Un couple acide/base est formé par un acide et sa base conjuguée. Par exemple :
\begin{itemize}
    \item \ce{CH3COOH/CH3COO-}
    \item \ce{NH4+/NH3}
\end{itemize}
En solution, les protons \ce{H+} ne restent jamais libres : ils sont toujours hydratés sous forme d'ion hydronium \ce{H3O+}.

\section{Autoprotolyse de l'Eau}
L'eau peut se comporter comme un acide ou une base. C'est ce qu'on appelle l'autoprotolyse de l'eau. Les demi-équations associées sont :
\begin{align*}
\ce{H2O &<=> H+ + OH-} \quad \text{(comportement basique)} \\
\ce{H2O + H+ &<=> H3O+} \quad \text{(comportement acide)}
\end{align*}
L'équation globale de l'autoprotolyse de l'eau est :
\[
\ce{2 H2O <=> H3O+ + OH-}
\]

\section{Produit Ionique de l'Eau}
Le produit ionique de l'eau, noté \(K_e\), est défini par :
\[
K_e = [\ce{H3O+}] \cdot [\ce{OH-}] = 1.0 \times 10^{-14} \text{ à } 25^\circ\text{C}.
\]

\section{pH et Échelle de pH}
\subsection{Définition}
Le pH est une mesure de l'acidité ou de la basicité d'une solution. Il est défini par :
\[
\text{pH} = -\log[\ce{H3O+}]
\]

\subsection{Échelle de pH}
\begin{center}
\begin{tabular}{|c|c|c|}
\hline
pH & Nature & Exemples \\ \hline
1 & Hyper Acide & \ce{HCl} dilué 100 fois \\ \hline
7 & Neutre & Eau pure à 25°C \\ \hline
14 & Hyper Basique & Hydroxyde de sodium \\ \hline
\end{tabular}
\end{center}

\section{Acides et Bases Forts et Faibles}
\subsection{Acides et Bases Forts}
\begin{itemize}
    \item Réagissent totalement avec l'eau.
    \item Exemples : \ce{HCl}, \ce{NaOH}.
\end{itemize}

\subsection{Acides et Bases Faibles}
\begin{itemize}
    \item Réagissent partiellement avec l'eau.
    \item Exemples : Acide acétique (\ce{CH3COOH}), Ammoniac (\ce{NH3}).
\end{itemize}

\section{Constante d'Acidité (Ka) et pKa}
\subsection{Constante d'Acidité}
La constante d'acidité, notée \(K_a\), est définie par :
\[
K_a = \frac{[\ce{A-}][\ce{H3O+}]}{[\ce{AH}]}
\]

\subsection{pKa}
Le pKa est défini par :
\[
\text{pKa} = -\log(K_a)
\]

\section{Coefficient de Dissociation}
Le coefficient de dissociation, noté \(\alpha\), est le rapport entre la quantité d'acide sous forme dissociée et la concentration initiale.

\section{Indicateurs Colorés}
Les indicateurs colorés changent de couleur en fonction du pH. Ils sont utilisés pour déterminer le pH d'une solution.

\section{Exercices}
\begin{enumerate}
\item Calculer le pH d'une solution d'acide chlorhydrique de concentration \(1.0 \times 10^{-2}\) mol/L.
\item Compléter le tableau suivant pour différentes concentrations d'acide acétique.
\end{enumerate}

\section{Conclusion}
Les réactions acido-basiques sont essentielles en chimie. Elles interviennent dans de nombreux processus naturels et industriels.

\end{document}
